\par\bigskip\Needspace{5\baselineskip}
\begin{center}\bfseries 2.\end{center}
\noindent\hspace*{\parindent}Zahlentabellen, wie die zuvor genannten, dienten damals wie später \textsc{Gauss} als numerisches Beobachtungsmaterial, um auf induktivem Wege arithmetische Gesetzmässigkeiten daraus zu erkennen, wie er denn in einem Briefe an den Hofrat \textsc{Zimmermann} vom 12. März 1797\footnote{Dieser Brief ist abgedruckt W. X\textsubscript{1}, S. 19—21.} sie als »vielleicht nützlich« bezeichnet, »in diesem Felde noch neue Erfindungen zu machen«. So geschah es denn auch, wie \textsc{Gauss} in der Vorrede der D. A. erzählt, dass zufällig bei einer anderen Arbeit im Anfange d. J. 1795 ein ausgezeichneter arithmetischer Satz (er meint den Satz des art. 108 vom quadratischen Charakter von \(-1\)) ihm aufstiess, der ihn nun reizte, seine tieferen Gründe und seinen Beweis aus denselben zu finden, und hieraus entsprangen dann schnell \textsc{Gauss}’ hierher gehörige theoretische Forschungen und Ergebnisse. Noch in demselben Jahre (1795) und zwar im Monat März (W. I, S. 476, Anm. zu art. 131) entdeckte \textsc{Gauss} durch Induktion das berühmte theorema fundamentale von den quadratischen Resten. Im Herbst d. J. siedelte er nach Göttingen über und fand hier schon am 8. April 1796 (Nr. 2 T.) seinen ersten Beweis des Fundamentaltheorems. Ohne Zweifel hat \textsc{Gauss} nun in Göttingen die bis dahin ihm fehlende Bekanntschaft (s. W. I, S. 6) mit den Arbeiten seiner Vorgänger \textsc{Euler}, \textsc{Lagrange}, \textsc{Legendre} schnell sich verschafft, von der seine Zitate in den D. A. Zeugnis geben. Nach seiner Aussage in der Vorrede zu den D. A. aber hatte er den grössten Teil ihrer ersten vier Abschnitte vorher schon »absolviert«. Hiernach und auf Grund der Daten, welche \textsc{Gauss} in seinem \textit{Tagebuch} für die auf jene Abschnitte bezüglichen Entdeckungen angegeben hat und die fast sämtlich in\newpage
\SourceLeaf{13}\ArithmeticHead{6}{L}{BACHMANN, ÜBER GAUSS’ ZAHLENTHEORETISCHE ARBEITEN.}
\noindent das Jahr 1796 fallen, wie der schon erwähnte Beweis des Fundamentaltheorems, mit dem er den Höhepunkt dieser Teile der D. A. erreichte, darf man annehmen, dass er den Inhalt der ersten vier Abschnitte schon 1796 besass. Ob er sie aber damals bereits niedergeschrieben hatte?

In seinem Nachlasse befindet sich ein Manuskript mit dem Titel \textit{Analysis Residuorum}, das als eine frühere Niederschrift der letzten Abschnitte der D. A. zu betrachten ist. Sie besteht aus drei Kapiteln, dem 6., 7. und 8. Das letztere kann nicht vor dem September 1797 entstanden sein, da nach den Angaben des \textsc{Gauss}schen \textit{Tagebuchs} die Auffindung seiner Sätze oder Beweise zum grössten Teil über 1796 hinaus bis zum 9. September 1797 reicht. In art. 367 desselben (W. II, S. 235) aber heisst es von der kubischen Gleichung für die drei \(\frac{\nu-1}{3}\)-gliedrigen Perioden aus \(\nu^{\mathrm{ten}}\) Einheitswurzeln: quam in Cap. VI a priori determinandam docuimus. Diese Gleichung ist (Nr. 39 T.) am 1. Oktober 1796 auf induktivem Wege von \textsc{Gauss} ermittelt, aber erst am 20. Juli 1797 (Nr. 67 T.) wirklich hergeleitet worden. Da nun das 6. Kapitel der A. R. genau mit einem unvollendeten Versuche, dies zu leisten, abbricht, jene Gleichung selbst noch nicht darin auftritt, darf man vermuten, dass dies Kapitel vor dem letztbezeichneten Datum verfasst sei. Was es sonst von der Kreisteilung enthält, geht nicht über das Jahr 1796 hinaus. Das 7. Kapitel ferner, das bei einer nicht auf den Text bezüglichen Randzeichnung das Datum »1796 Januar« trägt, enthält im Wesentlichen genau dasselbe, was den sechsten Abschnitt der D. A. bis zu der Stelle ausmacht, wo die Theorie der quadratischen Formen zur Anwendung kommt, und man möchte also nach dem, was über die darauf bezüglichen Tabellen gesagt ist, annehmen, dass es spätestens schon 1796 niedergeschrieben sei. Für das 6. Kapitel würde dann dasselbe gelten. Müsste man aber die Notiz Nr. 68 T. auf die in der Annotatio des art. 251 (W. II, S. 209) bemerkte Schwierigkeit beziehen, die dann erst am 21. Juli 1797 behoben worden wäre, so könnte das 6. Kapitel nicht vor diesem Tage verfasst sein. Dasselbe würde gelten, wenn man die Worte Error \textsc{Lagrangii}, welche in der Handschrift dem 7. Kapitel vorangestellt sind, auf \textsc{Lagranges} von \textsc{Gauss} (W. II, S. 249) getadelte Meinung beziehen müsste. Jedenfalls aber wird es richtig sein, wenn man die Niederschrift der A. R. auf 1796 bis 1797 datiert.
\clearpage\SourceLeaf{14}\ArithmeticHead{7}{R}{DIE ENTSTEHUNG DER D. A. DIE ANALYSIS RESIDUORUM.}
\noindent\hspace*{\parindent}Die Niederschrift der vier ersten Abschnitte der D. A. wird vor jener und in der Hauptsache vor 1797 geschehen sein, doch kann die endgültige Fassung der D. A. ihnen frühestens in diesem Jahre gegeben sein, da sich in \textsc{Gauss}’ \textit{Tagebuche} bis in letzteres hinein noch Angaben über gelungene Beweise für Sätze der vier ersten Abschnitte finden. In der Tat spricht \textsc{Gauss} in seinem Briefe an den Hofrat \textsc{Zimmermann} vom 12. März 1797\footnote{Siehe W. X\textsubscript{1}, S. 19.} von seinem Werke noch erst wie von einem Entwurfe, und seine Briefe an denselben vom 20. November und 24. Dezember 1797 zeigen, dass er erst in dieser Zeit beschäftigt ist, die ersten Bogen seines Werks für den Druck fertig zu stellen\begingroup\renewcommand{\thefootnote}{**)}\footnote{Auszüge aus den beiden letzteren Briefen sind abgedruckt in der Schrift: \textit{Karl Friedrich Gauss. Zwölf Kapitel aus seinem Leben} von L. \textsc{Hänselmann}. Leipzig 1878, S. 34—37. Die in dem Briefe vom 22. November (siehe a. a. O. S. 35) erwähnten Stellen accedere possunt S. 5 und eine andere S. 7 stehen W. I, S. 11 Zeile 9 v. u. resp. S. 13 Zeile 7—8.}\endgroup. Der fünfte Abschnitt kann nicht wohl vor dem 22. Juni 1796 (nach Nr. 15 T.) begonnen sein, hatte aber im Sommer 1798 schon eine dritte Darstellung erhalten (Brief von \textsc{Gauss} an \textsc{Bolyai} vom 29. Nov. 1798\begingroup\renewcommand{\thefootnote}{***)}\footnote{\textit{Briefwechsel zwischen C. F. Gauss und W. Bolyai}, herausgegeben von F. \textsc{Schmidt} und P. \textsc{Stäckel}. Leipzig 1899, S. 11.}\endgroup). Seine vierte Bearbeitung stammt diesem Briefe und den Notizen in \textsc{Gauss}’ \textit{Tagebuch} zufolge bestimmt aus dem Winter 1798/9. \textsc{Gauss} schreibt u. a., dass es ihm bei jeder folgenden Bearbeitung geglückt sei, die Sache auf eine solche Art auszuführen, dass es seine bei der vorhergehenden gehegten kühnsten Hoffnungen überstieg, und er werde das in ein paar Tagen zum vierten Male vollendet haben, was er im ganzen vorigen Sommer zum dritten Male ausarbeitete. Diesmal war es offenbar die Theorie der ternären quadratischen Formen, welche solches Gelingen herbeiführte und deren Studium er in jenen Tagen begann (s. Nr. 95 und 96 T.). Die vierte Bearbeitung des fünften Abschnitts wird in der Hauptsache dessen endgültige Fassung gewesen sein, wenngleich \textsc{Gauss} einen sehr wesentlichen Punkt der genannten Theorie erst im Februar 1800 erledigte (Nr. 103 T.). Es heisst weiter in dem erwähnten Briefe: »der sechste [Abschnitt] ist von keinem grossen Umfange; der siebente (der die Theorie der Polygone enthält) etwas grösser aber im Wesentlichen schon fertig und nur der letzte wird mich noch eine beträchtliche Zeit beschäftigen, da er die schwersten Materien enthält«. Dieser Teil ist bei Veröffentlichung der D. A.\newpage
\SourceLeaf{15}\ArithmeticHead{8}{L}{BACHMANN, ÜBER GAUSS’ ZAHLENTHEORETISCHE ARBEITEN.}
\noindent wegen zu grossen Umfangs des Werkes — jedenfalls auch, weil \textsc{Gauss} seinen Inhalt noch nicht ausgereift genug fand — unterdrückt worden.

Das Fragment, welches unter dem Titel Sectio octava in W. II, S. 510 veröffentlicht ist, wird als eine zur Fortsetzung der D. A. unternommene Neubearbeitung der A. R. anzusehen und erst nach Vollendung der D. A. zu setzen sein, da die Artikelnummern unmittelbar an die der D. A. sich anschliessen. Nachdem \textsc{Gauss} auch diese einstweilen bei Seite gelegt, nahm er erst 1808 in dem andern Fragmente (W. II, S. 243) seine Absicht wieder auf, jedoch hat er deren Ausführung über seinen andern gleichzeitigen Untersuchungen auch bald wieder fallen lassen.
